\section{Related work}\label{sec:related}

\paragraph{Unlearning methods and their evaluation.} Approximate unlearning for language
models is dominated by gradient ascent on the forget set, negative preference
optimisation \citep{zhang2024npo} and its reference-free variant SimNPO
\citep{fan2025simnpo}, and representation misdirection \citep{li2024wmdp}. They are
evaluated on benchmarks such as TOFU \citep{maini2024tofu} and MUSE \citep{shi2024muse}
whose forget-quality and leakage metrics are fixed-sample statistics compared against a
retrained model. A growing literature shows that these metrics overstate forgetting:
forget-quality metrics track surface behaviour rather than retained knowledge
\citep{wang2025towards}, output matching hides degenerate models \citep{yuan2025closer},
benign relearning restores erased content \citep{hu2025jogging}, relearning attacks can be
formalised and partly resisted \citep{fan2025relearning}, quantisation brings forgotten
knowledge back \citep{zhang2025quantization}, lightweight attacks recover unlearned
capabilities \citep{lucki2024adversarial}, and current evaluations can be made to
overstate or understate success \citep{feng2025existing}; see
\citet{liu2025rethinking} and \citet{xue2026survey} for surveys. Closest to our partial-unlearning experiment, \citet{li2026beliefs}
show that gradient-based unlearning shifts probability mass onto paraphrases of the
targeted answer, which they call squeezing, and propose an objective that accounts for
the model's own beliefs. We study the measurement side of the same phenomenon: whether a
certificate on the trained form can be misled, and what an audit needs to see it.

\paragraph{Membership inference as an audit.} Membership-inference attacks are the usual
instrument of unlearning audits. Score-based attacks threshold the loss, its
zlib-normalised version \citep{carlini2021extracting}, or the mean of the least likely
tokens, Min-K\% Prob \citep{shi2024mink}; reference-based attacks calibrate the loss by a
reference model, and LiRA \citep{carlini2022lira} calibrates per example with shadow
models. \citet{hayes2024inexact} show that unlearning evaluations based on
population-level attacks overstate privacy and that per-example calibration with
shadow models is needed. \citet{zhang2024mia} establish that membership-inference
scores cannot prove that a model was trained on given data unless the inclusion of that
data was randomised, the principle behind VOUCH's fair coins. For pre-trained language
models, \citet{duan2024membership} find that membership-inference attacks perform near
chance when members and non-members are drawn from the same distribution, and
\citet{das2024blind} show that classifiers that never query the model match or beat published
attacks on benchmarks whose members and non-members differ in distribution. Our matched- and
unmatched-holdout baselines (Section~\ref{sec:exp-baselines}) show what happens to a
membership-inference audit when that principle is not respected, and
Theorem~\ref{thm:detect} states the same confound in the language of a certificate.

\paragraph{Verification of unlearning.} \citet{thudi2022necessity} argue that
approximate unlearning cannot be audited from weights alone and that auditable
definitions are a prerequisite for any guarantee; \citet{zhang2024verification} show
that a dishonest provider can defeat existing families of verification strategies while
passing them; cryptographic proof-of-unlearning \citep{eisenhofer2025verifiable} attests
that the committed computation ran, which is complementary to whether forgetting holds
statistically. VOUCH \citep{dang2026vouch} made statistical verification anytime-valid
with paired canaries and betting e-processes
\citep{waudbysmith2024,ramdas2023game,shafer2021testing,grunwald2024safe}, building on
the secret-sharer methodology \citep{carlini2019secret}. Privacy auditing with
canaries \citep{steinke2023privacy} shares the planted design but estimates a lower
bound on the privacy parameter of a training algorithm, not the outcome of one
unlearning run. Copyright traps \citep{meeus2024traps} plant sequences in documents so
that their inclusion in a model's training data can be tested later, a planted design
closer in spirit to ours, without a paired control.
\citet{yu2025impossibility} show that retraining-based definitions of unlearning are not
well posed for multi-stage pipelines, one reason VOUCH avoids a retrained reference.
The special-casing provider that VOUCH's threat model excludes is the subject of
Section~\ref{sec:organic}: we do not remove the assumption, but we price it.

\paragraph{What is new here.} Relative to VOUCH, this paper (i) moves the canaries to
natural language and proves the exact null survives, (ii) states and proves the transfer
from canaries to organic records as a function of canary detectability, and measures
detectability for several designs including one built to make it zero, (iii) tests the
certificate against partial unlearning that suppresses the trained form, with two model
sizes and a union audit, (iv) compares with membership-inference audits on the same
runs, and (v) bounds the certification time through wealth growth rather than through
refutation.
